\begin{description}
  \item[学习单元（learning unit）] 围绕一个可观察学习产出组织的最小自学闭环。
  \item[独立 oracle（independent oracle）] 不由被测实现计算得出的期望结果。
  \item[命题（proposition）] 具有确定真假值、需要证明或反驳的陈述。
  \item[必要条件（necessary condition）] 结论成立时必须满足的条件。
  \item[充分条件（sufficient condition）] 一旦满足即可推出结论的条件。
  \item[反例（counterexample）] 满足命题假设但使结论失败的合法对象。
  \item[量纲（dimension）] 数量所属的单位类别，用于判断运算是否有意义。
  \item[度量空间（metric space）] 配备满足分离、对称和三角不等式之距离函数的集合。
  \item[Cauchy 序列（Cauchy sequence）] 尾部任意两项都能任意接近的序列。
  \item[完备性（completeness）] 每个 Cauchy 序列都收敛到空间内一点的性质。
  \item[一致收敛（uniform convergence）] 同一索引阈值同时控制定义域所有点误差的收敛。
  \item[压缩映射（contraction mapping）] 以严格小于 1 的统一因子缩短任意两点距离的映射。
  \item[$\sigma$ 代数（sigma algebra）] 包含全集并对补集和可数并封闭的事件集合族。
  \item[测度（measure）] 满足可数可加性的非负集合函数。
  \item[可测函数（measurable function）] 目标可测集的逆像仍是可测事件的函数。
  \item[简单函数（simple function）] 只取有限多个值的可测函数。
  \item[Lebesgue 积分（Lebesgue integral）] 由简单函数从下逼近构造的积分。
  \item[几乎处处（almost everywhere）] 除一个零测集外性质均成立。
  \item[$L^p$ 空间（Lp space）] $p$ 次绝对可积函数按几乎处处相等取等价类得到的完备范数空间。
  \item[本质上确界（essential supremum）] 忽略零测集后成立的最小几乎处处上界。
  \item[Hölder 不等式（Holder inequality）] 用共轭 $L^p$ 范数控制乘积积分的不等式。
  \item[Minkowski 不等式（Minkowski inequality）] 保证 $L^p$ 范数满足三角不等式的积分不等式。
  \item[Tonelli 定理（Tonelli theorem）] 对非负可测函数交换积分次序的定理。
  \item[Fubini 定理（Fubini theorem）] 在绝对可积条件下交换积分次序的定理。
  \item[乘积测度（product measure）] 从可测矩形上的测度乘积扩张得到的测度。
  \item[多元可微（multivariable differentiability）] 函数增量具有统一线性近似且相对余项消失。
  \item[梯度（gradient）] 表示标量函数一阶线性近似的偏导向量。
  \item[Jacobian] 表示向量值函数一阶线性近似的偏导矩阵。
  \item[Hessian] 描述标量函数局部二阶曲率的方阵。
  \item[矩阵微分（matrix differential）] 用线性微分和形状账本推导导数的记法。
  \item[正交投影（orthogonal projection）] 使残差位于目标子空间正交补的线性映射。
  \item[最小二乘（least squares）] 最小化残差平方和的估计方法。
  \item[谱分解（spectral decomposition）] 实对称矩阵的正交特征分解。
  \item[奇异值分解（singular value decomposition）] 任意矩阵的左右正交方向与奇异值分解。
  \item[条件数（condition number）] 问题对输入相对扰动的最坏敏感度。
  \item[概率空间（probability space）] 样本空间、事件 $\sigma$ 代数与概率测度组成的三元组。
  \item[随机变量（random variable）] 从概率空间到目标可测空间的可测映射。
  \item[推前分布（pushforward distribution）] 随机变量把底层概率测度推到目标空间得到的分布。
  \item[联合分布（joint distribution）] 随机向量在乘积空间上的分布。
  \item[边缘分布（marginal distribution）] 对联合分布的其余坐标积分或求和得到的分布。
  \item[重尾分布（heavy-tailed distribution）] 尾部衰减较慢并可能使某些矩不存在的分布。
  \item[条件期望（conditional expectation）] 给定子 $\sigma$ 代数时保持其事件上积分的可测随机变量。
  \item[塔式性质（tower property）] 对嵌套信息依次条件化等于直接对较粗信息条件化。
  \item[独立性（independence）] 联合概率对所有可测集合分解为边缘概率乘积。
  \item[条件独立（conditional independence）] 给定指定信息后联合条件分布分解为条件边缘乘积。
  \item[概率核（probability kernel）] 对状态可测、对事件为概率测度的条件分布映射。
  \item[几乎处处收敛（almost sure convergence）] 除零概率状态外逐点收敛。
  \item[依概率收敛（convergence in probability）] 偏差超过任意正阈值的概率趋零。
  \item[$L^p$ 收敛（Lp convergence）] 差的 $p$ 次绝对矩趋零。
  \item[分布收敛（convergence in distribution）] 分布函数在极限分布连续点收敛。
  \item[大数定律（law of large numbers）] 在明确矩与依赖条件下样本平均趋近总体均值。
  \item[中心极限定理（central limit theorem）] 标准化和在正则条件下趋近正态分布。
  \item[集中不等式（concentration inequality）] 控制有限样本偏离概率的非渐近上界。
  \item[估计量（estimator）] 把样本映射为参数估计的随机规则。
  \item[一致性（consistency）] 估计量依概率收敛到目标参数的性质。
  \item[渐近分布（asymptotic distribution）] 估计误差适当中心化与缩放后的极限分布。
  \item[异方差稳健标准误（heteroskedasticity-robust standard error）] 用三明治协方差估计且不要求条件误差方差恒定的标准误。
  \item[族错误率（family-wise error rate）] 一族真零假设中至少错误拒绝一个的概率。
  \item[错误发现率（false discovery rate）] 拒绝集合中错误发现比例的期望控制目标。
  \item[随机过程（stochastic process）] 同一概率空间上由索引参数化的随机变量族。
  \item[滤过（filtration）] 随时间递增、表示可用信息增长的 $\sigma$ 代数族。
  \item[Markov 性（Markov property）] 给定当前状态后未来条件分布不再依赖更早历史。
  \item[鞅（martingale）] 下一时点条件期望等于当前值的可积适应过程。
  \item[Poisson 过程（Poisson process）] 具有独立平稳 Poisson 增量的计数过程。
  \item[Brownian 运动（Brownian motion）] 具有连续路径及独立平稳高斯增量的过程。
  \item[二次变差（quadratic variation）] 沿细分割对路径增量平方求和的极限。
  \item[桥接模块（bridge module）] 根据诊断结果精确指回上册先修单元的短模块。
  \item[Capstone] 由干净复现、独立基线、失败案例和限制声明组成的综合能力证据。
\end{description}
